\chapter{CONCLUSION AND RECOMMENDATIONS}

This chapter summarizes the findings of the study and provides conclusions based on the empirical data. It also offers recommendations for future research directions.

%========================================================================
\section{Conclusions}
%========================================================================
Based on the experiments and analysis conducted, the following conclusions are drawn:

\begin{enumerate}
    \item The proposed Multiple Embedding Steganography (MES) framework successfully demonstrates that a parallel embedding architecture can significantly enhance the embedding capacity of RGB images compared to traditional sequential methods.
    \item The method achieved an average embedding capacity greater than 3.0 bits per pixel (bpp) across standard test images, validating its capability to conceal large payloads, including full-size color images.
    \item Despite the high payload, the visual quality of the stego-images remained high, with Peak Signal-to-Noise Ratio (PSNR) values consistently exceeding the 30 dB threshold. This confirms that the adaptive edge detection mechanism effectively masks the distortion in high-frequency regions.
    \item The parallel approach eliminates the inter-channel dependency found in sequential methods, allowing for more robust and independent data management across the Red, Green, and Blue channels.
\end{enumerate}

%========================================================================
\section{Recommendations}
%========================================================================
To further enhance the capabilities and robustness of the proposed framework, the following recommendations are suggested for future research:

\begin{enumerate}
    \item \textbf{Robustness Testing:} Future studies should evaluate the method's resilience against active attacks, such as image compression (JPEG), geometric transformations (rotation, scaling), and noise addition.
    \item \textbf{Frequency Domain Implementation:} Exploring the application of the parallel embedding concept in the frequency domain (e.g., DCT, DWT) could potentially offer better robustness against signal processing attacks.
    \item \textbf{Security Enhancement:} Integrating a cryptographic layer (e.g., AES encryption) before embedding would add an essential layer of security, ensuring that the hidden data remains confidential even if extracted.
    \item \textbf{Optimization:} Further optimization of the edge detection threshold using machine learning or genetic algorithms could potentially yield an even better trade-off between capacity and imperceptibility. Recent work treating the Canny hysteresis thresholds as explicit optimisation variables reports measurable gains in the capacity--fidelity trade-off and offers a concrete starting point \citep{singh2026edgeopt}.
    \item \textbf{Colour-Aware Steganalysis:} The security of a parallel colour architecture should be argued against detectors that model inter-channel dependence, not only against channel-agnostic ones. Independent modification of separate colour planes perturbs the natural correlation between those planes, and that perturbation is a documented detection cue \citep{amrutha2023rgbsteganalysis}. Evaluating MES against colour rich model features and modern convolutional detectors---using a panel rather than a single detector, since detector performance is itself unstable across conditions \citep{alrusaini2025robustness}---would convert the present imperceptibility evidence into a security claim.
    \item \textbf{Standardised Benchmarking:} The four classical test images support comparability with prior work but are too few to sustain a statistical security argument. Fragmented benchmarking is a recognised structural weakness of the data-hiding literature \citep{janok2026survey}, and evaluating on a standardised corpus such as StegBench \citep{andone2026stegbench} alongside a larger colour dataset would place future results on a firmer basis.
    \item \textbf{Correlation-Preserving Allocation:} Building on the previous point, a productive direction is to make inter-channel correlation preservation an explicit objective of the allocation rule rather than a constraint checked after the fact---for example, by distributing the payload so that the induced distortion across the three planes matches the cover's own channel statistics. This would address the principal security exposure of the parallel architecture without surrendering the capacity advantage that motivates it.
\end{enumerate}

%========================================================================
\section{Limitations}
%========================================================================
Three limitations bound the conclusions stated above and are recorded here so that the scope of the claims is not overstated. First, the evaluation set comprises four classical $512 \times 512$ images; this is adequate for characterising the capacity--quality trade-off across texture densities but not for a statistical claim about detectability. Second, the study is confined to passive-warden conditions in the spatial domain: no claim is made regarding payload survival under compression, resampling, or geometric transformation. Third, the comparison against learning-based image-in-image hiding is necessarily indirect, since those methods operate at payloads roughly two orders of magnitude lower and recover the secret approximately rather than bit-exactly \citep{janok2026survey}; the exact-recovery guarantee of the proposed spatial construction is an advantage for the target applications, but it means the two families are not measured on a common scale.
